% SOURCE_AUTHORITY: sga5_fr_workpass.tex, lines 6771--6792 (LF-normalized unit).
% SOURCE_SHA256: 791F4EFFC5E02832D5D77ED03518C8156D6F07E4C8238B03545DB93D883FBB28
\subsubsection*{2.2. Sistemas projetivos AR-nulos}
\begin{statement}{Definição 2.2.1}
Seja $X$ um objeto da categoria $P$. Diz-se que $X$ é AR-nulo se existe um inteiro $r\ge0$ tal que o morfismo canônico $X[r]\to X$ seja nulo.
\end{statement}
Em particular, é claro que um sistema projetivo AR-nulo verifica a condição (MLAR).

\begin{statement}{Proposição 2.2.2}
Seja
\[
0\longrightarrow X\longrightarrow Y\longrightarrow Z\longrightarrow0
\]
uma sequência exata de $P$. As seguintes afirmações são equivalentes:
\begin{enumerate}[label=(\roman*)]
\item $Y$ é AR-nulo;
\item $X$ e $Z$ são AR-nulos.
\end{enumerate}
Em outras palavras, a subcategoria plena $P_0$ de $P$, cujos objetos são os sistemas projetivos AR-nulos, é espessa em $P$.
\end{statement}

\begin{prooftext}[Demonstração]
É claro que (i) implica (ii). Reciprocamente, sejam $r$ e $t$ inteiros que satisfaçam a condição da definição 2.2.1 para $X$ e $Z$, respectivamente. Vê-se imediatamente que o inteiro $r+t$ satisfaz as condições da definição 2.2.1 para $Y$.
\end{prooftext}
